\section{Discussion}
STEI-ZStat provides EIS performance comparable to PalmSens4 over the validated 0.2~Hz--200~kHz range. At US\$25.36 per unit, it is the least expensive system listed in Table~\ref{tab:comparison} under the stated cost bases, although the margin over HunStat2 \cite{vamos2025hunstat2} is only about 8\% against its recalculated cost basis, and HunStat2 supports additional voltammetric methods. Its compact footprint and frequency range are also within the range of previously reported systems. Thus, the main contributions of this work are not the hardware cost, size, or frequency range, but the quantitative validation procedure used to assess measurement performance.

Measured repeatability was comparable between the instruments despite their different integration strategies (Table~\ref{tab:repeatability}). The longer STEI-ZStat integration windows at low frequencies may have contributed to this result, so the precision comparison should be interpreted in light of the differences in acquisition settings. Similarly, the lower median analytical magnitude deviation of STEI-ZStat does not indicate uniformly lower measurement error, because its high-frequency phase error was substantially larger.

Both estimates fall outside the component's $\pm$5\% tolerance but in the same direction, which is consistent with the aging of the Class II X7R dielectric after reflow. A maximum aging rate of 3\% per decade-hour is specified for X7R dielectrics \cite{kemet2025cots}, which, over the 134-day interval between receipt of the assembled dummy cell and measurement, corresponds to up to approximately 6\% downward drift. However, no aging rate was available for the specific component, so this explanation cannot be verified from the present data. The DC bias was 0~V, so DC-bias effects were absent. The 21~mV$_{rms}$ excitation was also well below the 1~V$_{rms}$ condition at which the capacitance was specified, but the effect of excitation amplitude on the actual capacitance was not characterized here.

Because the capacitor could not be independently measured in the assembled dummy cell, its nominal 33~nF value is the least certain component of the analytical reference. The 0.69~nF difference between the instruments' fitted capacitances, 2.3\% of the STEI-ZStat value, therefore remains unresolved and is reported rather than attributed to either instrument.

The main STEI-ZStat limitation is high-frequency phase error, which reaches 1.215$^\circ$ at 200~kHz compared with 0.329$^\circ$ for PalmSens4. Such behavior is consistent with the sensitivity of high-frequency EIS to parasitic impedances in the measurement path. The three assembled boards reproduced this deviation within 0.04$^\circ$ at 200~kHz and reproduced the shape of the deviation curve across the upper decade, indicating a systematic error rather than an isolated defect in one board. An error that is systematic in this sense is, in principle, correctable by calibration or compensation, although the responsible mechanism was not isolated here. The 200~kHz boundary should therefore be interpreted cautiously, particularly for measurements whose conclusions depend on small phase differences.

The validation separates statistical significance from practical equivalence. All fitted-parameter differences were statistically significant, but every 95\% CI remained within the predefined acceptance window. This distinction is important because significance testing can identify small systematic differences without determining whether they are practically relevant. The present approach quantifies both the difference and its uncertainty, and then determines whether the 95\% CI remains within a predefined acceptance window.

The validated range covers much of the frequency region commonly used for biosensing, including frequencies up to approximately 100~kHz reported in previous applications. However, the present validation begins at 0.2~Hz rather than the 0.01--0.1~Hz reported in some applications; because the AD5941 is specified to 0.015~Hz, this shortfall is a limitation of the validation rather than of the hardware. The validation also used 30~mV$_p$ rather than the 5--15~mV$_p$ excitation typical of small-signal biosensing, so the results do not directly establish performance at lower excitation amplitudes. Only a passive Randles cell with an impedance range of 560~$\Omega$--10.56~k$\Omega$, with a single time constant and 0~V DC bias, was evaluated. The results therefore establish instrument equivalence for the tested passive impedance range, not equivalent performance with functionalized electrodes, faradaic systems, or complete biosensors.

The primary comparison used a single STEI-ZStat unit. The three-board comparison bounds board-to-board variation, but those boards shared a single microcontroller module and a single dummy cell, and were measured by one operator, so the result does not establish the reproducibility of independently built complete instruments. All three boards also deviated from the nominal values in the same direction and by a similar amount, so the comparison bounds board-to-board scatter but not a common-mode scale error. The 1\% $R_{CAL}$ tolerance could still produce an absolute magnitude offset larger than the cross-instrument differences reported here. Replicates were acquired in separate blocks for each instrument rather than interleaved, and the measurement environment was not controlled or recorded; temporal and environmental effects therefore cannot be separated from instrument effects. Finally, redox measurements, functionalized electrodes, and the additional electroanalytical methods implemented in firmware remain unvalidated.